\section{Asking the customer}\label{sec:queries}

\begin{definition}[Interface]
The customer interface $K$ offers four fixed questions about a call to $a$: $\mathsf{provide}(p)$,
a typed value that must fit $p$'s format; $\mathsf{choose}(p, V)$, one of the options $V$, each
shown with its least trusted source; $\mathsf{confirm}(p, v)$, which accepts, rejects or replaces
$v$; and $\mathsf{approve}(o)$, a yes or no on the whole operation. The customer answers
truthfully, according to the operation they want. Answers are own-voice by construction and
cannot be quoted, and the adversary cannot answer them (Assumption~\ref{as:adv}(c)).
\end{definition}

Call a parameter of $a$ \emph{determined} by a memory if every history in its fibre gives it the
same value with $\Perm = 1$, and \emph{undetermined} otherwise.

\begin{proposition}[What must be asked]\label{prop:queries}
Let a system be sound and execute the operation the customer wants whenever one is wanted, and
let $k$ parameters be undetermined. Then, in the worst case over the fibre, it must obtain a value
for each of the $k$ from the customer: $k$ answers, or one form with $k$ fields. If it is to be
sound independently of how customer text is read (Theorem~\ref{thm:nomodel}), it also needs
$\mathsf{approve}$, even when $k = 0$.
\end{proposition}
\begin{proof}
If an undetermined parameter is not asked, the fibre contains two histories that want different
values for it, and the system executes the same one in both. The approval is
Theorem~\ref{thm:nomodel}.
\end{proof}

$G$ asks one question per parameter that is not bound: missing ones are provided, ambiguous ones
chosen, blocked ones confirmed. A proposal that no record carries is confirmed like a blocked value,
with the agent named as its source, or offered as one more button when the customer has values of
their own. When any question was asked, it adds the approval. So it meets the
bound with $k$ answers plus one approval for $k \geq 1$, and asks nothing at $k = 0$, which is the
case where it relies on $Q$, $R$, $L$ and $E$ (Corollary~\ref{cor:residual}). Choosing among $r$
candidates costs one $\mathsf{choose}$ where yes/no confirmations need up to $r - 1$.

\begin{remark}
This section is the weakest of the set. The counting is immediate, and ``number of questions''
depends on the interface. It is likely to become a remark in the paper rather than a theorem. A
sharper version would weigh questions by the information the customer supplies, in bits
conditional on memory, and ask which interface minimises it.
\end{remark}
